problem:  
Let \( (M^m,g) \) be a closed smooth Riemannian manifold, and let \( (N,d,\mu) \) be a compact metric–measure space that is **stratified \(RCD(K,N_{\mathrm{eff}})\)** in the following sense: \(N\) decomposes as \(N = N_{\mathrm{reg}} \cup \Sigma\), where \(N_{\mathrm{reg}}\) is a smooth Riemannian manifold of dimension \(n\), and each connected component of the singular set \( \Sigma_k \subset \Sigma \) has a neighborhood isometric to a metric cone  
\[
C(L_k) = \bigl([0,\varepsilon)\times L_k,\, d\rho^2 + \rho^2 h_{L_k},\, \rho^{n_k-1}\,d\rho \, d\mu_{L_k}\bigr)
\]
whose link \( (L_k,h_{L_k},\mu_{L_k}) \) is a compact \(RCD(K_k,n_k-1)\) space with a positive spectral gap \( \lambda_1(L_k) > 0 \).  
Assume each cone \(C(L_k)\) itself is \(RCD(K_k',N_{\mathrm{eff}})\) for some \(K_k'\ge K > -\infty\), so the space \(N\) has globally bounded synthetic Ricci curvature and satisfies a stratified \(RCD\) condition.

For \(u_0 : M \to N_{\mathrm{reg}}\) of finite Korevaar–Schoen energy \(E(u_0)\), denote by  
\[
u_t : [0,\infty)\to L^2(M,N)
\]
the Ambrosio–Gigli–Savaré \(L^2\)-gradient flow of the Dirichlet energy \(E(u)\); assume such a flow exists (as known for \(RCD\) targets).

**Question:**  
Does there exist a quantitative constant \( \varepsilon_*(N) > 0 \), depending only on \(K\), \(N_{\mathrm{eff}}\), and the minimal spectral gap \(\min_k \lambda_1(L_k)\), such that the following *energy–repulsion phenomenon* holds?

If \(E(u_0) < \varepsilon_*(N)\) and \(u_0(M)\subset N_{\mathrm{reg}}\) lies in a homotopy class representable entirely in \(N_{\mathrm{reg}}\), then  
\[
u_t(M) \subset N_{\mathrm{reg}},\qquad \forall t>0.
\]
In other words, does a universal small–energy bound—governed by the synthetic lower curvature of the links—guarantee that the \(RCD\) harmonic map heat flow from a smooth domain never reaches the singular strata of the target?  
Equivalently, do the curvature–dimension bounds on the links induce enough convexity of the squared distance to the apex in the cone models to force analytic avoidance from the singularities?

Why is it a "good" problem:  
This question proposes a precise analytic–geometric bridge between **synthetic Ricci lower bounds** and **nonlinear parabolic regularization** in singular spaces.  It asks whether the curvature–dimension structure on the cone links, quantified by their spectral data, provides a parabolic barrier preventing the harmonic map heat flow from encountering the singular strata.  

The problem is simple to state—does small energy imply the flow avoids singularities?—yet pursuing it requires uniting diverse tools: metric–gradient–flow theory, variational calculus in \(RCD\) spaces, spectral geometry of links, and ε–regularity principles from geometric analysis.  Solving it would uncover whether synthetic Ricci positivity manifests dynamically as *curvature-induced repulsion* for energy flows, potentially forming the foundation for a generalized small–energy regularity theory of harmonic map heat flows into stratified \(RCD\) spaces.  Its simplicity of statement, cross-field depth, and potential to deepen the analytic understanding of synthetic curvature make it an excellent and forward-looking research-level problem.